\documentclass[11pt,a4paper]{article}

% ---------- packages ----------
\usepackage[utf8]{inputenc}
\usepackage[T1]{fontenc}
\usepackage{lmodern}
\usepackage{microtype}
\usepackage[margin=1in]{geometry}
\usepackage{amsmath,amssymb}
\usepackage{enumitem}
\usepackage{titlesec}
\usepackage{xcolor}
\usepackage{hyperref}
\hypersetup{
  colorlinks=true,
  linkcolor=black,
  citecolor=black,
  urlcolor=blue!55!black,
  pdftitle={AI Applications in Research: Notes on Claude Code},
  pdfauthor={Scheduled routine summary}
}
\usepackage{parskip}

\titleformat{\section}{\normalfont\Large\bfseries}{\thesection}{0.7em}{}
\titleformat{\subsection}{\normalfont\large\bfseries}{\thesubsection}{0.6em}{}

% ---------- metadata ----------
\title{AI Applications in Research \\
\large A short survey, with emphasis on Claude Code}
\author{Compiled from public sources by a scheduled routine \\
        \small carusomatteo17@gmail.com}
\date{2026-08-20}

\begin{document}
\maketitle

\begin{abstract}
This note surveys how large language model (LLM) agents --- and in
particular Anthropic's \emph{Claude Code} --- are being folded into
empirical research workflows in 2026. It draws on Scott Cunningham's
ongoing ``Claude Code'' series on \emph{Scott's Mixtape} Substack, the
\texttt{MixtapeTools} repository, recent reproducibility benchmarks in
the social sciences, and adjacent write-ups from academic economists.
The intent is a scannable overview rather than an exhaustive review:
what the tools are, which workflow patterns have emerged, what the
evidence says about reliability, and where the open problems sit.
\end{abstract}

\tableofcontents

% ==============================================================
\section{Introduction}
% ==============================================================

Two shifts have converged over 2024--2026. First, general-purpose LLMs
crossed a capability threshold at which they can reliably read a
repository, execute code, install dependencies, and iterate against a
concrete objective. Second, \emph{coding agents} --- terminal-native
harnesses around those models --- exposed that capability in a form
researchers already know how to use: a shell, a file tree, a git
history. Claude Code, Codex, Cursor's agent mode, and Aider are the
most visible instances. Empirical researchers have started to adopt
them not as autocomplete, but as a new layer in the research stack,
sitting between the researcher and the raw data.

This document focuses on Claude Code because (a) it currently leads
independent reproducibility benchmarks in the social sciences and
(b) it has the best-documented academic workflow to date, thanks to
Scott Cunningham's public series.

% ==============================================================
\section{What Claude Code is}
% ==============================================================

Claude Code is Anthropic's official command-line agent. Practically,
it is a terminal program that (i) accepts a natural-language goal,
(ii) reads and writes files in a working directory, (iii) runs shell
commands and inspects their output, and (iv) loops until the goal is
met or the user intervenes. Around that core it exposes several
extensibility points relevant to research:

\begin{itemize}[leftmargin=1.2em]
  \item \textbf{Skills} --- reusable, named prompt-and-tool bundles
        that ship with a repository or a user's dotfiles and load
        on demand (\texttt{/skill-name}).
  \item \textbf{Sub-agents} --- named agent types with restricted
        toolsets (e.g.\ a read-only \emph{Explore} agent, a
        \emph{code-reviewer} agent) that the main loop can delegate
        to in parallel.
  \item \textbf{Hooks} --- shell commands run on lifecycle events
        (\texttt{SessionStart}, \texttt{PostToolUse}, \ldots),
        useful for e.g.\ auto-installing project dependencies at
        session start on a fresh clone.
  \item \textbf{MCP servers} --- external tool providers exposed
        over the Model Context Protocol; commonly wired to GitHub,
        Google Drive, Gmail, Calendar, or a lab-specific data
        catalogue.
  \item \textbf{Web execution} --- a managed, ephemeral container
        variant that runs on Anthropic's infrastructure and can be
        triggered from GitHub or a schedule, so the same agent
        works interactively and unattended.
\end{itemize}

% ==============================================================
\section{Emerging workflow patterns}
% ==============================================================

Several patterns recur across published academic write-ups
(Cunningham; Goldsmith-Pinkham; Sant'Anna; Matray; Backman; Crawfurd).

\subsection{Claude as thinking partner, not autopilot}

The dominant posture in Cunningham's \texttt{MixtapeTools} is
\emph{iterative steering}: the researcher decomposes the task,
Claude executes when the next move is unambiguous, and surfaces a
question when it is not. Long autonomous runs are avoided for
substantive analytical work; they are reserved for mechanical passes
(formatting, migrations, cross-language ports).

\subsection{Cross-language replication audits}

A signature \texttt{MixtapeTools} pattern is the ``Referee~2'' or
``health inspector'' protocol: after a project appears complete, a
\emph{fresh} Claude instance --- with no memory of the build --- is
pointed at the artefacts and asked to replicate the pipeline
independently in \emph{different} languages (typically Python, R, and
Stata). Divergences between the three independent runs are the
finding. This inverts the traditional replication package: rather
than shipping code and hoping it runs elsewhere, the researcher
generates independent implementations up front.

\subsection{Presentation and writing pipelines}

Beamer decks, Quarto reports, and \LaTeX{} manuscripts are built and
maintained inside the same repo the analysis lives in, with the
agent handling the mechanical parts (figure regeneration, table
formatting, citation hygiene, cross-references). This is where the
productivity delta is largest and least controversial.

\subsection{Research management as data}

Cunningham exposes a ``research GTD'' layer --- hypotheses, open
questions, evidence, decisions --- as plain-text files under version
control that Claude can read and update. The agent becomes a
disciplined note-taker rather than a source of new claims.

% ==============================================================
\section{Evidence on reliability}
% ==============================================================

The most-cited benchmark to date is \emph{SocSci-Repro-Bench}
(Alizadeh et al., 2026), which evaluates coding agents on 221 tasks
drawn from 13 substantive social-science domains. Headline numbers:

\begin{itemize}[leftmargin=1.2em]
  \item Task-level accuracy: Claude Code \textbf{93.4\%}, Codex
        \textbf{62.1\%} (a 31.3~pp gap).
  \item Paper-level accuracy: Claude Code \textbf{78.0\%}, Codex
        \textbf{35.8\%} (a 42.2~pp gap).
  \item The authors argue the delta reflects \emph{agentic}
        capability (reading code, installing deps, debugging,
        executing) rather than memorised recall of results.
\end{itemize}

Complementary work on a many-analysts economics study finds that
agent outputs cluster near the human median with \emph{tighter}
dispersion and no extreme outliers --- suggestive that agents behave
as low-variance, moderately-skilled analysts rather than as one more
idiosyncratic team.

% ==============================================================
\section{Open problems}
% ==============================================================

Cunningham is explicit that Claude Code introduces new failure modes
even as it removes old ones. Recurring themes across the community:

\begin{itemize}[leftmargin=1.2em]
  \item \textbf{Silent overconfidence.} An agent that finishes
        without errors is not the same as an agent that finished
        correctly. Cross-language audits and adversarial verifiers
        exist precisely because a single run's ``done'' signal is
        weak.
  \item \textbf{Context loss and drift.} Long sessions accumulate
        stale assumptions; short, well-scoped sub-agents beat one
        long chat.
  \item \textbf{Reviewer legibility.} Diffs authored by an agent
        may be locally correct but stylistically alien; house-style
        skills and explicit review passes help.
  \item \textbf{Reproducibility of the agent itself.} Model
        upgrades, harness changes, and skill drift mean the ``same''
        prompt can produce different outputs six months later; pin
        model IDs and record harness versions in the paper's
        replication package.
  \item \textbf{Governance.} Data-use agreements, IRB scope, and
        vendor data-retention terms need to be reviewed before an
        agent is given raw subject data.
\end{itemize}

% ==============================================================
\section{Getting started}
% ==============================================================

A minimum viable setup for an empirical researcher looks roughly
like this:

\begin{enumerate}[leftmargin=1.4em]
  \item Install Claude Code and authenticate.
  \item Clone \texttt{MixtapeTools} and copy the skills you need
        into your project's \texttt{.claude/} directory.
  \item Add a \texttt{SessionStart} hook that (re)installs the
        project's dependencies, so a fresh clone --- or a fresh
        cloud session --- comes up green.
  \item Wire an MCP server for wherever your data lives (Drive,
        S3, a lab database).
  \item Adopt one review discipline (e.g.\ Referee~2 cross-language
        replication, or a \texttt{/code-review} pass gated on
        merge) so ``the agent said it was done'' is never the last
        word.
\end{enumerate}

% ==============================================================
\section{Selected sources}
% ==============================================================

\begin{itemize}[leftmargin=1.2em]
  \item Scott Cunningham, \emph{Online version of The Remix is now
        here}, Scott's Mixtape Substack, 2026-08-17.
        \url{https://causalinf.substack.com/p/online-version-of-the-remix-is-now}
  \item Scott Cunningham, \emph{Claude Code} series (index).
        \url{https://causalinf.substack.com/s/claude-code}
  \item \texttt{scunning1975/MixtapeTools} (GitHub).
        \url{https://github.com/scunning1975/MixtapeTools}
  \item M.\ Alizadeh et al., \emph{AI Coding Agents Can Reproduce
        Social Science Findings}, arXiv:2606.11447.
        \url{https://arxiv.org/abs/2606.11447}
  \item Anthropic, \emph{Coding agents in the social sciences}.
        \url{https://www.anthropic.com/research/coding-agents-social-sciences}
  \item P.\ Sant'Anna, \emph{Claude Code Academic Workflow --- \LaTeX,
        Quarto, Research Automation}.
        \url{https://psantanna.com/claude-code-my-workflow/}
  \item P.\ Goldsmith-Pinkham, \emph{Getting Started with Claude
        Code: A Researcher's Setup Guide}.
        \url{https://paulgp.substack.com/p/getting-started-with-claude-code}
  \item A.\ Matray, \emph{Custom Claude Code Skills for Research}.
        \url{https://adrienmatray-ai.com/skills.html}
  \item C.\ Backman, \emph{Claude Code in VS Code --- For Academic
        Economists}. \url{https://claesbackman.com/claude-code-guide.html}
  \item L.\ Crawfurd, \emph{Claude Code Skills for Academic
        Research}. \url{https://lcrawfurd.github.io/claude-skills/}
\end{itemize}

\vfill
\noindent\rule{\linewidth}{0.4pt}\\
{\small Compiled 2026-08-20 by a scheduled Claude Code routine.
Content is a synthesis of public sources; details attributed to
individual posts should be verified against the live articles before
citation. Substack domains were unreachable from the routine's
runtime environment, so post-level detail was reconstructed from
search snippets and adjacent public pages.}

\end{document}
